\documentclass[10pt,a4paper]{article}
\usepackage[margin=18mm]{geometry}
\usepackage{amsmath,booktabs,hyperref,pdfpages}
\hypersetup{colorlinks=true,urlcolor=blue,linkcolor=blue}
\setlength{\parindent}{0pt}\setlength{\parskip}{5pt}
\begin{document}
\begin{center}{\Large $\Lambda$--$^4$He high-$k^*$ systematics}\\[4pt]
Au+Au fixed target, 3.85 GeV, 2021 P24iy \quad 11 November 2026 (JST)\end{center}
\textbf{Finding.} The high-$k^*$ decrease is reproduced by the new histogram production.
Tightening the nuclear TPC pull alone changes the central-sample ratio little.
The decrease depends strongly on nuclear rigidity and also on pair $m_T$.
This identifies useful conditioning variables; it does not establish the physical or detector origin.

\textbf{Definition.} All quoted values are ratios of counts integrated over $0.7\leq k^*<0.8$ GeV/$c$,
scaled by $M_N/S_N$ from $N=[0.40,0.50)$ GeV/$c$:
\[ C_W=\frac{S_W}{M_W}\frac{M_N}{S_N}. \]
They are raw $\Lambda$ signal-mass-gate CFs, without purity corrections.
The reference interval is a comparison convention; no physical plateau is established.
Each cut is applied separately to a saved two-dimensional histogram.
$R=|\mathbf p_{\rm nucleus}|/2$ is in GeV/$c$; $m_T$ is in GeV/$c^2$.

\begin{center}\begin{tabular}{lcc}\toprule
Selection & Centrality 0--20\% & Centrality 20--80\%\\\midrule
Inclusive & 0.9079 $\pm$ 0.0033 & 0.8942 $\pm$ 0.0083 \\
$|n\sigma|<0.6$ & 0.9093 $\pm$ 0.0049 & 0.9161 $\pm$ 0.0128 \\
$|n\sigma|<1.0$ & 0.9065 $\pm$ 0.0040 & 0.9011 $\pm$ 0.0102 \\
$R<1.0$ & 0.8576 $\pm$ 0.0051 & 0.8816 $\pm$ 0.0101 \\
$1.0\leq R<1.5$ & 0.9152 $\pm$ 0.0057 & 0.9128 $\pm$ 0.0170 \\
$1.5\leq R<2.0$ & 0.9699 $\pm$ 0.0071 & 0.9822 $\pm$ 0.0333 \\
$2.0\leq R<5.0$ & 1.0002 $\pm$ 0.0165 & 1.0545 $\pm$ 0.1320 \\
$2.4\leq m_T<2.6$ & 0.8899 $\pm$ 0.0052 & 0.8929 $\pm$ 0.0107 \\
$2.8\leq m_T<3.0$ & 0.9667 $\pm$ 0.0106 & 1.0277 $\pm$ 0.0559 \\
\bottomrule\end{tabular}\end{center}

For centrality 0--20\%, the $|n\sigma|<1$ minus inclusive difference is
$-0.00139\pm0.00222$ after accounting for nested-count overlap.
The $R<1$ minus inclusive difference is $-0.05026\pm0.00412$.
These are count-model errors, not calibrated event-level significance tests.
The inclusive 0--20\% ratio is $0.8651\pm0.0028$ in $0.8\leq k^*<1.0$ GeV/$c$.
The inclusive 0--20\% linear slope over $0.5\leq k^*<0.8$ is
$-0.3299\pm0.0148$ $(\mathrm{GeV}/c)^{-1}$, with the full normalization covariance.

\textbf{Interpretation limits.} The high-rigidity selection also requires TOF and a mass-squared window;
the low-rigidity selection does not require TOF. Their difference cannot isolate a TOF effect.
A tighter pull window can change the rigidity composition; the joint $R$--pull distribution is absent.
Rigidity-correlated contamination is therefore not excluded. No new physics cut is adopted here.
\newpage
\textbf{Species and mass controls.} The same high-$k^*$ decrease is weaker in the selected $^3$He control.
Changing its mass hypothesis does not reproduce the full $^4$He decrease.
The two $^3$He representations use the same candidates and are correlated.
\begin{center}\begin{tabular}{lcc}\toprule
Control & 0--20\% & 20--80\%\\\midrule
$\Lambda$--$^4$He & 0.9079 $\pm$ 0.0033 & 0.8942 $\pm$ 0.0083 \\
$\Lambda$--$^3$He & 0.9583 $\pm$ 0.0014 & 0.9467 $\pm$ 0.0051 \\
$^3$He with $^4$He mass & 0.9578 $\pm$ 0.0014 & 0.9419 $\pm$ 0.0053 \\
\bottomrule\end{tabular}\end{center}
The low and high halves of the $\Lambda$ signal window show similar decreases.
The sidebands are closer to unity; their different kinematics and composition prevent interpreting them as an independent purity correction.
\begin{center}\begin{tabular}{lcc}\toprule
$\Lambda$ mass region & 0--20\% & 20--80\%\\\midrule
Signal low half & 0.9051 $\pm$ 0.0043 & 0.9011 $\pm$ 0.0112 \\
Signal high half & 0.9115 $\pm$ 0.0050 & 0.8856 $\pm$ 0.0124 \\
Low-mass sideband & 0.9629 $\pm$ 0.0228 & 1.0183 $\pm$ 0.0548 \\
High-mass sideband & 0.9819 $\pm$ 0.0131 & 0.9674 $\pm$ 0.0264 \\
\bottomrule\end{tabular}\end{center}
\textbf{Production reproducibility.} With the same TH2D projections, the old/new inclusive 0--20\% values are
$0.90781/0.90788$; for 20--80\% they are $0.89431/0.89421$.
The productions contain strongly overlapping events; their difference is not tested as independent statistics.
Changing the normalization range changes the vertical scale and does not remove the declining shape.

\textbf{Input and validation.} New SUMS ID \texttt{1CE33994716F7896423A80AB994BAB14}.
The verified merge contains 26,945 of 26,994 subjobs, with 49 excluded; it is a partial merge.
It contains 2,099,784,951 input and 2,066,742,256 selected events.
ROOT 5.34/38 read 1,656 $k^*$ TH1/TH2 histograms from each production without changing either input.
The scan covers 1,680 series, 5 grouped plus 9 fine centralities, and 1,176 projection/partition checks.
Independent checks reproduce 318 numerical results and 24 finite-difference gradient checks.
Legacy TH2F/TH2D differences are confined to the $^3$He mixed-event $k^*\geq1$ overflow;
the maximum is 31 counts, $2.52\times10^{-7}$ relative, outside the plotted analysis range.

\textbf{Uncertainties.} Stored count variances are propagated through count integration, normalization reuse,
and overlap of nested cuts. Slopes use the full bin covariance from normalization.
The saved histograms do not recover event reuse, shared mixed-pair correlations, or model/purity systematics.
Missing-count bins are shown as unavailable, rather than zero.

\textbf{What this sample supports.} Independent pull, rigidity and $m_T$ projections, mass regions,
centrality and species controls are available.
Candidate storage and joint 8D Sparse were disabled for capacity reasons.
Simultaneous rigidity--pull cuts, low-rigidity TOF-tagged CFs, and arbitrary event-level remixing cannot be reconstructed from this file.
Pre-PID histograms retain track quality, species rigidity and TOF requirements and use the existing physical mixing pool;
they are diagnostic views, not an unbiased alternate production.

\textbf{Next discriminating measurement.} With candidates/PID held fixed, compare narrower rigidity or $m_T$ composition and
controlled mixing in multiplicity, vertex, run/field and event plane. Such a remix requires retained candidate records or a separate PicoDst rerun.
No new production is submitted in this study.

\textbf{Atlas guide.} Following pages: species (1--2), mass regions (3--4), pull (5--7), rigidity (8--10),
$m_T$ (11--12), pre-PID (13--14), normalization (15--16), and old/new production (17--18).
All configurations, counts and statistical checks are retained beside this report.
\includepdf[pages=-,landscape=true,pagebox=cropbox]{cf_highk_systematics_20261111.pdf}
\end{document}
